\section{Formal coefficients and analytic traces}
\label{chap:formal-analytic-determinants}

A formal logarithm and a holomorphic logarithm satisfy the same
coefficient identities, but they answer different questions. A formal
logarithm determines each coefficient by finitely many operations.
A holomorphic logarithm also has a domain of convergence and a choice
of branch. Continuing that branch can require a new surface even when
the original coefficients are completely determined.

The prime determinant makes this distinction explicit. Its coefficients
belong to a natural Banach algebra precisely in its initial
regularization half-plane. The continued germ then lives on the surfaces
constructed in Sections~\ref{sec:prime-continuation} and~\ref{sec:finite-covers}. The same
diagonal matrices act on a formal product and on a Hilbert space, but
their traces have different continuity properties in those two settings.
We construct these comparisons, beginning with one Gaussian variable.

\subsection{Specializing coefficients without evaluating a formal variable}
\label{sec:da-specialization}

Let $R$ be a commutative algebra over $\mathbb Q$ and let $t$ be a
commuting formal variable. The ring $R[[t]]$ consists of sequences
$\sum_{n\geq0}a_nt^n$ with $a_n\in R$. The coefficient of $t^n$ in
a product is $\sum_{j=0}^na_jb_{n-j}$. Thus multiplication uses a finite
sum at every degree. The $t$-adic topology declares two series close
when their coefficients agree through a prescribed degree.

For $u\in tR[[t]]$, set
\[
 \exp(u)=\sum_{r\geq0}\frac{u^r}{r!},\qquad
 \log(1+u)=\sum_{r\geq1}\frac{(-1)^{r+1}}r u^r.
\]
At degree $n$, only $r\leq n$ can contribute. These are therefore
formal operations with no condition on the size of $a_n$.
Their usual inverse identities hold coefficientwise. For completeness,
write $E(z)=\sum z^r/r!$ and $L(z)=\sum_{r\geq1}(-1)^{r+1}z^r/r$.
Formal differentiation gives $E'=E$ and $L'=(1+z)^{-1}$.
The product and chain rules give
\[
 (L(E(z)-1))'=1,\qquad
 \bigl(E(L(z))(1+z)^{-1}\bigr)'=0.
\]
Their constant terms are zero and one. Characteristic zero therefore
gives $L(E(z)-1)=z$ and $E(L(z))=1+z$.
Substitution of $u\in tR[[t]]$ preserves these finite coefficient
identities.

\begin{proposition}[Coefficient specialization and its limit]
\label{prop:da-specialization}
A unital algebra homomorphism $\phi:R\to\mathbb C$ extends coefficientwise
to an algebra homomorphism
\[
 \phi[[t]]:R[[t]]\longrightarrow\mathbb C[[t]].
\]
It commutes with the displayed formal exponential and logarithm.
For any $t_0\in\mathbb C^\times$, there is no unital
$\mathbb C$-algebra homomorphism $\mathbb C[[t]]\to\mathbb C$
that sends $t$ to $t_0$.
\end{proposition}
\begin{proof}
Each product, exponential, or logarithm coefficient is a finite
polynomial expression over $\mathbb Q$ in the input coefficients.
Applying $\phi$ to that expression proves the first assertions.
For the last assertion, the series
\[
 (1-t/t_0)^{-1}=\sum_{n\geq0}(t/t_0)^n
\]
is an inverse in $\mathbb C[[t]]$. The proposed homomorphism would
send $1-t/t_0$ to zero and its product with this inverse to one.
This is impossible.
\end{proof}

For the commutative scalar sector of formal Gaussian elimination, take
$A=\mathbb C[x_1,\ldots,x_d]$. Its coefficient algebra is
$A[\hbar,\hbar^{-1}][[t]]$. Each $t$-coefficient is polynomial in the
fields and a finite Laurent polynomial in $\hbar$. A scalar
specialization of $A$, together with a number $h_0\ne0$, therefore
defines a coefficientwise map sending $\hbar$ to $h_0$.
The lower Laurent bound may depend on the $t$-degree; no uniform bound
is needed for this map. Its output is still a formal series in $t$.
The proposition supplies no evaluation at a nonzero numerical $t$.
If odd field generators are present, a scalar algebra map sends them
to zero. The same coefficient argument applies to the even interaction.

This statement differs from specializing an arbitrary Laurent series
in $\hbar$. In $\mathbb C((\hbar))$, the element $\hbar-h_0$ is
invertible for every $h_0\ne0$, since its constant term is nonzero.
The same unit argument excludes a unital $\mathbb C$-algebra map
$\mathbb C((\hbar))\to\mathbb C$ with $\hbar\mapsto h_0$.
The coefficient ring is part of the comparison.

\subsection{A convergent Gaussian logarithm and a divergent interaction}
\label{sec:da-gaussian}

Fix a real number $b>0$. Let $\langle f\rangle_b$ be the normalized
Gaussian integral of a polynomial $f$ with density
$e^{-x^2/(2b)}$ on $\mathbb R$. Integration by parts gives
\[
 \langle x^{2r}\rangle_b=(2r-1)!!\,b^r,\qquad
 \langle x^{2r+1}\rangle_b=0,
 \qquad (-1)!!=1.
\]
Indeed, the integral of the derivative of
$x^{2r-1}e^{-x^2/(2b)}$ is zero and gives the recurrence.
Every polynomial times this Gaussian tends to zero at infinity.
For polynomials, the same functional is
\[
 \langle f\rangle_b=
 \left.\exp\!\left(\frac b2\partial_x^2\right)f\right|_{x=0}.
\]
Both sides have the displayed values on every monomial, and the
operator series terminates on each polynomial.
For a coefficient parameter $b$, the polynomial operator
$T_b=\exp(b\partial_x^2/2)$ is called Gaussian contraction; each
polynomial involves only finitely many derivatives. Extend it
coefficientwise to $\mathbb C[x,b][\hbar,\hbar^{-1}][[t]]$.
The numerical integral above is its specialization at a real $b>0$,
followed by evaluation at $x=0$. This is an equality of linear
functionals on polynomials, not of algebra homomorphisms.

For a real polynomial $J\in\mathbb R[x]$, define the formal interaction weight
\[
 Z_J(t)=\sum_{n\geq0}\frac{\langle J^n\rangle_b}{n!}t^n,
 \qquad W_J(t)=\log Z_J(t)\in t\mathbb R[[t]].
\]
This uses the formal operations above and does not assume that the
integral of $e^{tJ}$ converges.
For $I\in t\mathbb C[x,b][\hbar,\hbar^{-1}][[t]]$, define the
formal effective interaction by
$\mathcal W_b(I)=\hbar\log(T_b(e^{I/\hbar}))$.
The coefficientwise exponential and logarithm are defined because
$I/\hbar$ has positive $t$-order and $T_b(1)=1$.
After fixing $\hbar=h_0\ne0$, taking $I=h_0tJ$, and evaluating
at $x=0$, the weight is $Z_J$ and $\mathcal W_b(I)/h_0$ is $W_J$.

\begin{proposition}[Quadratic and quartic interactions]
\label{prop:da-gaussian}
For $J=x^2/2$, the formal series is the Taylor series of
\[
 Z_J(t)=(1-bt)^{-1/2},\qquad
 W_J(t)=-\tfrac12\log(1-bt)
\]
on $|t|<1/b$, with $Z_J(0)=1$ and $W_J(0)=0$.
The normalized integral defines $Z_J$ on $\operatorname{Re}t<1/b$.
Its continuation around $t=1/b$ changes sign and has exact local
ramification index two.

For $J=x^4$, the formal series $Z_J$ has radius of convergence zero.
The normalized real integral of $e^{tx^4}$ diverges for $t>0$ and
is finite for $t\leq0$.
\end{proposition}
\begin{proof}
For the quadratic interaction, the exponent in the integral is
$-(1-bt)x^2/(2b)$. On a compact subset of
$\operatorname{Re}t<1/b$, this integrand and each of its derivatives
in $t$ are bounded by an integrable polynomial times a fixed Gaussian.
Differentiation under the integral is therefore valid.
For real $t<1/b$, a real change of variable in the normalized
Gaussian integral gives $(1-bt)^{-1/2}$.
Both functions are holomorphic on this half-plane, so the equality
holds there. The logarithm is the branch with value zero at $t=0$;
$1-bt$ has positive real part throughout this domain.

For $|t|<1/b$, the series of absolute values of the integrand is
$e^{|t|x^2/2}$. It is Gaussian-integrable. Termwise integration gives
exactly the formal coefficients. The singularity at $t=1/b$ makes
the Taylor radius exactly $1/b$.
The logarithmic derivative is
\[
 \frac{Z_J'}{Z_J}=\frac{b}{2(1-bt)}
                  =-\frac1{2(t-1/b)}.
\]
Its integral about a small positively oriented loop is $-\pi i$,
so the multiplier is $-1$. On the local coordinate
$t=1/b+u^2$, one has $Z_J=c/u$ with $c^2=-1/b$.
This is a simple pole on the two-sheeted local surface. An unramified
meromorphic function has an integral order and cannot have the
base exponent $-1/2$, so two is minimal.

For the quartic interaction, the coefficient of $t^n$ is
$a_n=(4n-1)!!\,b^{2n}/n!$. Its consecutive ratio is
\[
 \frac{a_{n+1}}{a_n}
 =\frac{(4n+3)(4n+1)}{n+1}b^2\longrightarrow\infty.
\]
For every nonzero complex $t$, the terms $a_nt^n$ therefore fail
to tend to zero. This proves the zero radius.
For real $t>0$, the exponent $tx^4-x^2/(2b)$ tends to positive
infinity, so the integral diverges. For $t\leq0$, its integrand is
bounded by the original Gaussian. The integral is then finite.
\end{proof}

The quartic logarithm remains an exact formal series. It cannot be
the Taylor series of a holomorphic germ whose exponential is $Z_J$,
because exponentiating such a germ would give a convergent Taylor
series for $Z_J$. Thus a connected formal expansion does not by itself
produce a holomorphic effective interaction. The quadratic example
also shows that convergence near the initial point does not remove
fractional monodromy elsewhere.

\subsection{The coefficient algebra of a prime product}
\label{sec:da-dirichlet}

The variables in an Euler product are related by multiplication of
integers. To retain these relations before numerical evaluation, let
\[
 \mathscr D_{\mathrm{form}}=
 \left\{\sum_{n\geq1}a_n[n]:a_n\in\mathbb C\right\},
 \qquad [r][s]=[rs].
\]
The brackets denote formal basis symbols, not values of a function.
Multiplication is Dirichlet convolution:
\[
 (a*b)_n=\sum_{d\mid n}a_d b_{n/d}.
\]
There are only finitely many divisors of $n$, so every coefficient
exists. The unit is $[1]$, and $a\mapsto a_1$ is an augmentation.
Let $F_N$ consist of the series with $a_n=0$ for $n<N$.
These subspaces give the coefficient topology. Their intersection
is zero, compatible finite coefficient lists give a unique series,
and $F_NF_M\subset F_{NM}$.
In particular, if $a_1=0$, the coefficient of $[n]$ in $a^r$ is
zero for $2^r>n$. Exponential and logarithm consequently exist on
$F_2$ and $1+F_2$, respectively, by the same finite coefficient
formulas as before.

For a positive integer $m$, put
\[
 \mathscr L_m=-\sum_p\sum_{j\geq m}\frac{[p^j]}j,
 \qquad \mathscr D_m=\exp(\mathscr L_m).
\]
Each coefficient of the first sum is finite. In a fixed quotient
by $F_N$, only primes and exponents with $p^j<N$ occur. Applying
the exponential identity for a finite sum in this quotient gives
\begin{equation}
 \mathscr D_m=\prod_p
 \left((1-[p])\exp\!\sum_{j=1}^{m-1}\frac{[p^j]}j\right).
 \label{eq:da-formal-product}
\end{equation}
The product in~\eqref{eq:da-formal-product} is defined coefficientwise. Unique prime
factorization ensures that, for each $[n]$, only primes dividing
$n$ can contribute. The coefficient of $[p^m]$ in $\mathscr D_m$
is $-1/m$, since the corresponding single-prime logarithm begins
with $-[p^m]/m$.

The formal product does not allow arbitrary numerical substitution.
For any fixed $s\in\mathbb C$, the element
$1-2^s[2]$ is invertible in $\mathscr D_{\mathrm{form}}$, with
inverse $\sum_{r\geq0}2^{rs}[2^r]$.
A unital $\mathbb C$-algebra homomorphism sending $[n]$ to $n^{-s}$
would send this unit to zero. No such homomorphism exists on the
whole formal algebra. For example,
$\sum_{r\geq0}r![2^r]$ has a finite coefficient at every integer,
but its proposed numerical sum diverges at every $s$.

\subsection{A Banach algebra on which evaluation exists}
\label{sec:da-banach}

To control numerical evaluation, fix a real number $\sigma>0$ and
define
\[
 \mathscr D_\sigma=
 \left\{a:\|a\|_\sigma=
                 \sum_{n\geq1}|a_n|n^{-\sigma}<\infty\right\}.
\]
Here the coefficients have their usual complex absolute value.
This norm imposes a condition absent from the formal coefficient
topology. The underlying coefficient inclusion
$\mathscr D_\sigma\subset\mathscr D_{\mathrm{form}}$ preserves
addition and multiplication; it is not an assertion that the norm
topology equals the formal topology.

\begin{lemma}[Summable coefficients and holomorphic evaluation]
\label{lem:da-evaluation}
The space $\mathscr D_\sigma$ is a commutative unital Banach algebra.
On $\mathbb H_\sigma=\{s:\operatorname{Re}s>\sigma\}$, the map
\[
 a\longmapsto\operatorname{ev}(a)(s)=\sum_{n\geq1}a_nn^{-s}
\]
is an algebra homomorphism to bounded holomorphic functions, with
$\sup_{\mathbb H_\sigma}|\operatorname{ev}(a)|\leq\|a\|_\sigma$.
It is the unique norm-continuous extension of evaluation on finite
Dirichlet sums. Moreover,
$\operatorname{ev}(a)(s)\to a_1$ as real $s\to+\infty$.
\end{lemma}
\begin{proof}
Absolute summation and $(rs)^{-\sigma}=r^{-\sigma}s^{-\sigma}$ give
\[
 \|a*b\|_\sigma
 \leq\sum_{r,s\geq1}|a_r||b_s|(rs)^{-\sigma}
 =\|a\|_\sigma\|b\|_\sigma.
\]
Multiplication is therefore continuous and $\|[1]\|_\sigma=1$.
For a norm-Cauchy sequence, every coefficient is Cauchy. Call its
coefficientwise limit $a$. Bounds on finite sums pass to the limit,
and then increasing the finite set shows that $a$ is summable.
The same argument applied to differences proves norm convergence
to $a$. This proves completeness.

The evaluation series converges absolutely and uniformly on
$\operatorname{Re}s\geq\sigma$. The product of two such series can
be rearranged absolutely, giving the convolution formula.
On $\operatorname{Re}s\geq\sigma+\delta$, where $\delta>0$,
its derivative series is bounded in absolute sum by
\[
 \sum_n |a_n|(\log n)n^{-\sigma-\delta}
 \leq\left(\sup_{x\geq1}(\log x)x^{-\delta}\right)\|a\|_\sigma.
\]
The supremum is finite. Uniform convergence of the function and
derivative series on each such half-plane permits termwise
differentiation, proving holomorphy. Truncations converge in norm,
which proves the uniqueness assertion. Finally, dominated
convergence in the absolutely summable majorant
$|a_n|n^{-\sigma}$ gives the limit $a_1$.
\end{proof}

The prime-harmonic divergence proved in
Theorem~\ref{thm:euler-fredholm} makes the following convergence
threshold sharp, including its boundary.

\begin{theorem}[Exact coefficient domain of the regularized determinant]
\label{thm:da-coefficient-domain}
For every integer $m\geq1$ and real $\sigma>0$,
\[
 \mathscr L_m\in\mathscr D_\sigma
 \quad\Longleftrightarrow\quad
 \mathscr D_m\in\mathscr D_\sigma
 \quad\Longleftrightarrow\quad m\sigma>1.
\]
On $\operatorname{Re}s>1/m$, their evaluations are exactly
\[
 \operatorname{ev}(\mathscr L_m)(s)=L_m(s),\qquad
 \operatorname{ev}(\mathscr D_m)(s)=\Delta_m(s),
\]
with the logarithm, real prime powers and normalization of
Theorem~\ref{thm:prime-regularization} and
Section~\ref{sec:finite-covers}. In particular,
$\Delta_m(s)\to1$ along the positive real axis.
\end{theorem}
\begin{proof}
If $m\sigma>1$, then
\[
 \|\mathscr L_m\|_\sigma
 =\sum_p\sum_{j\geq m}\frac{p^{-j\sigma}}j
 \leq\frac1{m(1-2^{-\sigma})}\sum_p p^{-m\sigma}
 <\infty.
\]
The last sum is bounded by $\sum_{n\geq2}n^{-m\sigma}$, which
converges by comparison with its integral.
In the Banach algebra, the exponential series converges in norm,
with norm at most $\exp\|\mathscr L_m\|_\sigma$.
Its coefficient at $[n]$ equals the formal exponential coefficient,
since only finitely many powers contribute. Thus its sum is
$\mathscr D_m$.

If $m\sigma\leq1$, the coefficients at $[p^m]$ of both
$\mathscr L_m$ and $\mathscr D_m$ have absolute value $1/m$.
Consequently each proposed norm is at least
\[
 \frac1m\sum_p p^{-m\sigma}\geq\frac1m\sum_p\frac1p=\infty.
\]
This proves both necessary conditions, including the boundary.

Given $\operatorname{Re}s>1/m$, choose
$1/m<\sigma<\operatorname{Re}s$.
The evaluation homomorphism of Lemma~\ref{lem:da-evaluation}
commutes with the norm-convergent exponential series. Hence
\[
 \operatorname{ev}(\mathscr D_m)(s)
 =\exp\left(-\sum_p\sum_{j\geq m}\frac{p^{-js}}j\right).
\]
The prime powers use the real logarithm of each prime, exactly as
in the initial determinant definition. The right side is
$\exp L_m(s)=\Delta_m(s)$ on this half-plane.
The coefficient at $[1]$ of $\mathscr L_m$ is zero and that of
$\mathscr D_m$ is one. The limiting normalization follows from
Lemma~\ref{lem:da-evaluation}.
\end{proof}

The theorem identifies the actual germ used for continuation.
For $m\geq2$, continuing it does not extend evaluation on the whole
formal algebra. It also does not place its original coefficients
in a larger weighted algebra across $m\sigma=1$.
For instance, at $m=2$ the prime-square coefficients already force
divergence at $\sigma=1/2$, while the continued determinant has
local exponent $1/2$ at $s=1/2$.
The continuation is a statement about that holomorphic germ on its
branched surface, with total residue and period character as
constructed in Sections~\ref{sec:prime-continuation}
and~\ref{sec:finite-covers}.

For $n>m$, the formal identity
\[
 \mathscr D_n=
 \exp\left(\sum_p\sum_{j=m}^{n-1}\frac{[p^j]}j\right)
 \mathscr D_m
\]
holds coefficientwise. Every factor belongs to $\mathscr D_\sigma$
when $\sigma>1/m$. Evaluation therefore gives exactly the multiplier
$H_{m,n}$ of Theorem~\ref{thm:order-comparison} on that overlap.
Outside it, the comparison uses the common continuation cover of
that theorem; the formal identity alone specifies no branch.

\subsection{The same diagonal matrix on two different spaces}
\label{sec:da-operators}

The formal recipient can be specified without an analytic norm.
Let $\mathbb C[[x,y]]/\mathbb C$ be the vector space of formal
functions modulo constants. Its Hamiltonian Lie bracket is
$\{f,g\}=\partial_xf\,\partial_yg-\partial_yf\,\partial_xg$
modulo constants. Both derivatives and every output coefficient are
well-defined in the coefficient topology. Fix a bijection between the primes and the nonconstant
monomials $x^ay^b/(a!b!)$. The normalization only rescales the
monomial basis. In these coordinates the formal space is
\[
 \widehat V=\mathbb C^{\{p\text{ prime}\}}
       \cong\mathbb C[[x,y]]/\mathbb C.
\]
Give each copy of $\mathbb C$ the discrete topology and give
$\widehat V$ the product topology. A continuous linear operator on
$\widehat V$ is exactly a row-finite matrix. Indeed, each output
coordinate must be a continuous scalar functional, so it depends
on finitely many input coordinates. Conversely, finite support of
each row makes every finite collection of output coordinates
continuous. The operator topology requires agreement of whole rows
in each prescribed finite collection.

The same finite sequences form a dense subspace of the Hilbert space
$V_2=\ell^2(\{p\})$, now over the usual complex field, with its
declared orthonormal basis $e_p$. For $s\in\mathbb C$, define the
diagonal matrix
\[
 K(s)e_p=p^{-s}e_p.
\]
It is a continuous endomorphism of $\widehat V$ for every $s$.
No summation over primes is required for a fixed output coordinate.

\begin{proposition}[Operator domains and scalar traces]
\label{prop:da-operator-domains}
Write $\sigma=\operatorname{Re}s$. On $V_2$, the diagonal matrix
$K(s)$ is bounded exactly when $\sigma\geq0$, and compact exactly
when $\sigma>0$. For every integer $m\geq1$,
\[
 K(s)\in \mathcal S_m(V_2)\quad\Longleftrightarrow\quad m\sigma>1.
\]
On that half-plane, $s\mapsto K(s)$ is holomorphic in the
$\mathcal S_m$ norm. For $j\geq m$, the ordinary Hilbert trace is
\[
 \operatorname{Tr}K(s)^j=\sum_p p^{-js},\qquad
 L_m(s)=-\sum_{j\geq m}\frac{\operatorname{Tr}K(s)^j}{j},
\]
where the final double series is absolutely and locally uniformly
convergent.
\end{proposition}
\begin{proof}
For $\sigma\ge0$ the diagonal entries have supremum $2^{-\sigma}$;
for $\sigma<0$ they are unbounded. Finite diagonal truncations converge
in operator norm exactly when $\sigma>0$. At $\sigma=0$ the images
of the coordinate unit vectors form an orthonormal sequence, proving
noncompactness. The singular values, Schatten threshold and
norm-holomorphy are those of Theorem~\ref{thm:euler-fredholm}.
The trace formulas and absolute local uniform convergence are the
normalized logarithm of Theorem~\ref{thm:prime-regularization}.
\end{proof}

For $\sigma<0$, the diagonal matrix still defines a densely defined
closed operator on $V_2$, with domain
\[
 \mathcal D(K(s))=
 \left\{v\in\ell^2:\sum_p p^{-2\sigma}|v_p|^2<\infty\right\}.
\]
Finite sequences make the domain dense. If $v^{(n)}\to v$ and
$K(s)v^{(n)}\to w$ in $\ell^2$, convergence in each coordinate gives
$w_p=p^{-s}v_p$. Since $w\in\ell^2$, $v$ lies in the displayed
domain and $K(s)v=w$. This proves closedness. It does not supply
a trace-class power in that region.

The formal and Hilbert realizations agree algebraically on finite
sequences and on every coordinate where both are defined.
Their topologies do not agree. The coordinate inclusion
$V_2\to\widehat V$ is not continuous: $n^{-1}e_2$ tends to zero
in Hilbert norm, while its $2$-coordinate never becomes zero in
the discrete field. Likewise, $s\mapsto K(s)$ is not continuous
in the formal operator topology on any open complex domain, since
its first diagonal entry $2^{-s}$ is not locally constant.
These facts prevent a silent change from a formal completion to
a holomorphic family of operators.

\subsection{Trace compatibility and Hamiltonian translation}
\label{sec:da-trace-comparison}

For a continuous finite-rank operator on $\widehat V$, write
$T=\sum_{j=1}^k v_j\otimes\lambda_j$, where each $\lambda_j$ depends
on finitely many coordinates, and define
\[
 \operatorname{tr}_{\mathrm{fin}}T=\sum_{j=1}^k\lambda_j(v_j).
\]
Such an expression exists: choose a basis of the image, then finitely
many coordinates separating its basis vectors. Inverting that finite
coordinate matrix expresses the image coordinates as continuous
functionals. On any finite-dimensional subspace containing the image,
the displayed sum is the ordinary matrix trace of the restriction.
Enlarging this subspace adds a zero quotient operator, so the value is
independent of all choices. For a continuous $B$ and finite-rank $T$,
both $BT$ and $TB$ have trace $\sum_j\lambda_j(Bv_j)$.
On a Hilbert space with inner product linear in the first argument,
the corresponding rank-one operator is $h\mapsto\langle h,g\rangle f$
and its trace is $\langle f,g\rangle$. These definitions agree on
their common finite-coordinate operators.

Let $\epsilon^p(v)=v_p$ and let $P_p=e_p\otimes\epsilon^p$ be the rank-one coordinate
projection. Its finite-rank trace is one in both realizations.
For a finite diagonal sum, both traces equal the sum of its
diagonal entries. This is the precise common finite calculation.

On summable diagonal sequences, the Hilbert trace extends to
\[
 \ell^1(\{p\})\longrightarrow\mathbb C,
 \qquad (a_p)\longmapsto\sum_p a_p,
 \qquad \left|\sum_p a_p\right|\leq\sum_p|a_p|.
\]
Truncations converge in this norm. For the formal operator topology,
the projections $P_p$ tend to zero as the prime index escapes every
finite set, but all their traces equal one. Hence finite-rank trace
has no continuous linear extension to the full formal operator
completion, even when the scalar target has its usual topology.
The Hilbert trace extends by a different convergence condition.

An algebraic cyclic extension is also impossible on all row-finite
operators. Enumerate the basis as $e_0,e_1,\ldots$ and define
$Le_0=0$, $Le_{n+1}=e_n$, $Re_n=e_{n+1}$.
Both matrices are row-finite and
\[
 LR=1,\qquad RL=1-P_0,\qquad [L,R]=P_0.
\]
Any cyclic scalar functional on that algebra must vanish on $P_0$,
contradicting the rank-one trace. On the span of
$e_0,\ldots,e_{N-1}$, the truncated commutator is instead
$P_0-P_{N-1}$. Both finite traces are one and cancel.
The second term vanishes in the formal operator limit, but its
trace does not. This identifies the precise discontinuity.

The prime labeling introduces no Hamiltonian symmetry by itself.
The following computation tests that issue on one translation.

\begin{proposition}[The prime diagonal is not translation-invariant]
\label{prop:da-translation}
For the chosen bijection of primes with normalized monomials,
let $q_n$ label $f_n=y^{n+1}/(n+1)!$, for $n\geq0$.
Let $a_x$ be the formal Hamiltonian action
$a_x f=\{x,f\}=\partial_y f$ modulo constants.
For every real $s>0$ and every $n\geq1$,
\[
 [a_x,K(s)]f_n
 =(q_n^{-s}-q_{n-1}^{-s})f_{n-1}\ne0.
\]
Thus $K(s)$ is not an endomorphism commuting with the full
Hamiltonian action on the formal space.
\end{proposition}
\begin{proof}
The factorial normalization gives $a_xf_n=f_{n-1}$ for $n\geq1$.
Applying the two compositions gives $q_n^{-s}f_{n-1}$ and
$q_{n-1}^{-s}f_{n-1}$. The primes $q_n$ and $q_{n-1}$ are
distinct, and the function $r\mapsto r^{-s}$ is strictly decreasing
on positive reals when $s>0$. Their difference is nonzero.
\end{proof}

The comparison also has an explicit realization on coefficient functions.
Introduce one commuting variable $u_p$ for each basis vector and put
\[
 A_{\mathrm{coeff}}=\prod_\alpha\mathbb C u^\alpha,
 \qquad O_v=\sum_p v_pu_p\quad(v\in\widehat V),
\]
where $\alpha$ runs over finitely supported families of nonnegative
integers. The product is well-defined because a fixed multi-index has
only finitely many decompositions $\alpha=\beta+\gamma$.
Use the product topology with discrete coefficients. The formulas
\[
 \mathcal E(u^\alpha)=\left(\sum_p\alpha_p\right)u^\alpha,
 \qquad
 \delta_{K(s)}(u^\alpha)=
       \left(\sum_p\alpha_pp^{-s}\right)u^\alpha
\]
define continuous derivations of this algebra. Each coefficient sum is
finite, and additivity of the displayed weights under
$\alpha=\beta+\gamma$ proves the product rule.
They satisfy
\[
 \mathcal E(1)=0,\qquad \mathcal E(O_v)=O_v,
 \qquad \delta_{K(s)}(O_v)=O_{K(s)v}.
\]
Thus the Euler derivation is the derivation associated with the identity
kernel, while the prime diagonal assigns distinct weights to the chosen
linear coordinates. These are actual maps on the same coefficient
algebra. Their construction uses no infinite scalar trace. The
translation commutator above proves that the prime assignment does not
identify an operator commuting with the Hamiltonian action.

We have constructed the coefficient evaluation map, identified its
exact domain for $\mathscr D_m$, and compared the resulting
determinant with actual trace-class powers on the prime Hilbert
space. The continued scalar germ retains the same normalization and
the same minimal surfaces and finite-cover criteria of
Sections~\ref{sec:prime-continuation} and~\ref{sec:finite-covers}, including the single-valued
case $m=1$. These constructions supply neither a cyclic trace on the
full formal operator algebra nor an equivariant arithmetic matter
representation. An identification with quantum current amplitudes
would require a specified representation and trace ideal, compatible
contractions, and convergence of the associated spatial integrals.
Those maps are additional mathematical data.
